\section{Dense vs. MoE 配置回顾与补充}

\subsection{$d_{\text{model}}$ 的本质}

$d_{\text{model}}$（即 \texttt{hidden\_size}）是贯穿整个模型始终的核心维度：
\begin{itemize}
  \item 每层的输入 token embedding 维度为 $d_{\text{model}}$
  \item 每层的输出 token embedding 维度仍为 $d_{\text{model}}$
  \item $d_{\text{model}}$ 的一致性保证了\textbf{残差连接}的可行性（变换前后维度相同，可以直接相加）
\end{itemize}

\subsection{中间维度与 FFN 结构}

\begin{table}[H]
\centering
\caption{FFN/MoE 中间维度对比}
\begin{tabular}{lcc}
\toprule
& \textbf{Dense (32B)} & \textbf{MoE (30B-A3B)} \\
\midrule
中间维度 $d_{\text{ff}}$ & 25600 ($5 \times d_{\text{model}}$) & 768（每个 Expert） \\
映射方向 & $5120 \to 25600$（up） & $2048 \to 768$（down!） \\
Expert 数 & --- & 128 \\
激活 Expert 数 & --- & 8 \\
\bottomrule
\end{tabular}
\end{table}

\begin{warningbox}{MoE Expert 的中间维并非 ``up''}
对于 Dense 模型，FFN 的中间维 $25600 > d_{\text{model}} = 5120$，是一个「升维」过程。但 MoE 的每个 Expert 中间维 $768 < d_{\text{model}} = 2048$，实际上是一个「降维」过程。虽然代码中仍沿用 $W_{\text{up}}$ 的命名，但本质上每个 Expert 是更「瘦」的。
\end{warningbox}

\subsection{Decoder Sparse Stack 配置}

Qwen3-30B-A3B 的配置中有两个关键参数：
\begin{itemize}
  \item \texttt{decoder\_sparse\_step}：控制多少层使用一次 Dense MLP（取模为 0 时用 Dense MLP）
  \item \texttt{mlp\_only\_layers}：指定哪些层使用标准 Dense MLP
\end{itemize}

在当前配置下，\textbf{所有 48 层的 FFN 部分都使用 MoE}，没有任何层回退到 Dense MLP。

\subsection{本章小结}

MoE 模型在维度上全面「缩窄」（$d_{\text{model}}$、头数、层数），通过 128 个独立 Expert 来补偿表达能力。每个 Expert 的 FFN 是「窄」的（768 维），但多个 Expert 组合后等效宽度可达 6144。

